\subsection{Projective programs under the assumption}
\label{rank:projective-program-input}

The argument for a projective base uses a terminating program to reach a nef
adjoint and, after decreasing the nef data, a program ending in a Mori fiber
space. The conditional input is termination in the pseudo-effective case.
We first state the exact program result and then derive that case from the
projective good models in Proposition~\ref{proj:good-models}.

\begin{proposition}[A conditional generalized program]
\label{proj:generalized-mmp}
Assume Assumption~\ref{asm:log-iitaka}. Let
\((V,B+\boldsymbol M)\) be a projective generalized dlt rational pair,
where \(V\) is \(\Q\)-factorial, \(B\geq0\), and the nef b-divisor
\(\boldsymbol M\) is determined by a nef rational Cartier divisor on a
projective birational model. Put \(D=K_V+B+M_V\).

There is a choice of \(D\)-MMP with ample scaling which terminates. If
\(D\) is pseudo-effective, its endpoint \(V_m\) has nef adjoint \(D_m\).
If \(D\) is not pseudo-effective, its endpoint has a Mori fiber contraction
to a projective variety of smaller dimension. The birational steps and their
endpoints stay \(\Q\)-factorial generalized dlt, and the forward birational
map extracts no divisors.
In the pseudo-effective case, on a common resolution one has an actual
rational divisor identity
\[
 u^*D=v^*D_m+F,\qquad F\geq0\text{ and }v\text{-exceptional}.
\]
\end{proposition}

For the pseudo-effective alternative, we first derive the relative weak
Zariski decompositions used by the minimal-model existence theorem.

For a projective morphism, first replace its image by its normal Stein
factor. We use the relative convention of \cite[Section~2]{TsakanikasXie}:
a divisor is pseudo-effective over the base when its restriction to a very
general fiber of this surjective morphism is pseudo-effective. An NQC weak Zariski decomposition of \(D\) over the
base is a birational equality \(h^*D\equiv_Z P+N\), where \(P\) is a
nonnegative real combination of relatively nef rational Cartier divisors
and \(N\geq0\). The decomposition below has rational Cartier parts.

\begin{lemma}[Relative canonical decomposition]\label{proj:relative-wzd}
Assume Assumption~\ref{asm:log-iitaka}. Let \(V\) be a smooth
quasi-projective complex variety and let \(f:V\to Z\) be projective, with
\(Z\) normal and quasi-projective. If \(K_V\) is pseudo-effective over
\(Z\), then on a normal variety \(Y\) with a projective birational morphism \(h:Y\to V\) there
are rational Cartier divisors \(P,N\) such that
\[
 h^*K_V\equiv_Z P+N,\qquad P\text{ nef over }Z,\qquad N\geq0.
\]
\end{lemma}

\begin{proof}
We use the construction in the proof of
\cite[Proposition~3.3]{External037}. Its absolute input is a nef-plus-
effective rational decomposition for the canonical divisor of a smooth
projective variety with pseudo-effective canonical class. This is supplied
here by Proposition~\ref{proj:good-models} and its common-resolution
comparison.

Replace the image of \(f\) by its normal Stein factor. This does not change
relative numerical classes or contracted curves, and now \(f\) is
surjective with connected fibers. Its very general fiber is smooth
projective with pseudo-effective canonical divisor, hence is non-uniruled
by \cite[Theorem~0.2 and Corollary~0.3]{BDPP}. Compactify the projective
morphism and resolve away from \(V\). We obtain
\(\bar f:\bar V\to\bar Z\), with \(\bar V\) smooth projective and
\(\bar Z\) normal projective, unchanged over \(Z\). If \(\dim Z=0\),
\(V\) is already projective and the absolute input proves the lemma.

Suppose \(\dim Z>0\), and let \(r:\widetilde Z\to\bar Z\) be a smooth
projective resolution. Choose a sufficiently positive very ample divisor
\(A_Z\) on \(\bar Z\) and a general \(B\in|2A_Z|\). The pulled-back systems
on both \(\widetilde Z\) and \(\bar V\) are base point free. Bertini makes
\(r^*B\) and \(\bar f^*B\) smooth nonempty reduced divisors. Their double
covers
\[
 \pi_Z:Z^\sharp\to\widetilde Z,
 \qquad \pi:V^\sharp\to\bar V
\]
are smooth integral projective varieties. The canonical formulas are
\[
 K_{Z^\sharp}\sim_{\Q}\pi_Z^*(K_{\widetilde Z}+r^*A_Z),
 \qquad
 K_{V^\sharp}\sim_{\Q}\pi^*(K_{\bar V}+\bar f^*A_Z).
\]
The first line is big when \(A_Z\) is sufficiently positive. The rational
map \(V^\sharp\dashrightarrow Z^\sharp\) has the same very general fibers
as \(f\). A covering family of rational curves on \(V^\sharp\) would
either dominate a covering family of rational curves on \(Z^\sharp\), or
cover its very general fibers by vertical rational curves. Both are
impossible. Thus \(V^\sharp\) is non-uniruled and \(K_{V^\sharp}\) is
pseudo-effective by BDPP.

Apply the absolute decomposition to \(V^\sharp\). Moving the rational
principal difference in the canonical formula into the nef part gives an
actual rational divisor decomposition for a pullback of
\(\pi^*(K_{\bar V}+\bar f^*A_Z)\). Take a common equivariant resolution of
that model and its conjugate under the covering involution. On this smooth
projective resolution \(g:T\to V^\sharp\), average the two decompositions.
We obtain
\[
 g^*\pi^*(K_{\bar V}+\bar f^*A_Z)=\bar P+\bar N,
 \qquad \bar P\text{ nef},\quad \bar N\geq0,
\]
where \(\bar P\) and \(\bar N\) are invariant rational Cartier divisors.

Let \(q:T\to\bar Y\) be the finite quotient by the involution. The invariant
composite \(\pi g\) induces a projective birational morphism
\(\bar h:\bar Y\to\bar V\), with \(\bar Y\) normal and projective. An invariant
rational divisor descends by dividing its coefficient at an upstairs prime
by the ramification index. The descended divisor is rational Cartier when
the original divisor is: after clearing denominators, choose one Cartier
equation on a semilocal neighborhood of the finite orbit above a point
(a Cartier divisor is principal on a semilocal ring), and multiply all its
translates. The product is invariant and has divisor
equal to the group order times the original divisor. It descends to a local
equation for a multiple of the downstairs divisor. This argument includes
stabilizers and ramification primes.

Accordingly \(\bar P=q^*P\) and \(\bar N=q^*N\) for rational Cartier
divisors on \(\bar Y\). Lifting curves through the finite map shows that
\(P\) is nef, and coefficientwise descent shows that \(N\) is effective.
Finite pullback is injective on rational divisors, so the displayed identity
descends. Restrict it to \(Y=\bar h^{-1}(V)\). The term pulled
back from \(A_Z\) is numerically trivial over \(Z\), giving
\(h^*K_V\equiv_ZP+N\), as required.
\end{proof}

\begin{proof}[Proof of Proposition~\ref{proj:generalized-mmp}]
Lemma~\ref{proj:relative-wzd} and
\cite[Theorem~5.2]{TsakanikasXie}, applied to the generalized pair
\((U/Z,0+0)\) with \(U\) smooth, prove relative minimal-model existence for smooth varieties
by dimension induction. The hypothesis of that theorem is relative smooth
minimal-model existence one dimension lower; the lemma supplies its other
hypothesis, an NQC weak Zariski decomposition. The induction begins in
dimension zero. Then \cite[Theorem~5.4]{TsakanikasXie} supplies minimal
models for every pseudo-effective NQC generalized lc pair in the relevant
dimension. Their Theorem~2.7 gives existence in the Birkar--Shokurov sense.

The nef data in the statement are NQC, because a nef rational Cartier
divisor is a positive multiple of a nef Cartier divisor. In the
pseudo-effective case the preceding paragraph therefore gives the
Birkar--Shokurov model required by
\cite[Theorem~4.2]{TsakanikasXie}. In the other case, the non-pseudo-effective
alternative of that theorem applies directly. To meet its scaling
hypothesis, take an effective sufficiently positive ample rational divisor
\(A\) whose general components have sufficiently small coefficients. On
a fixed log resolution carrying \(\boldsymbol M\), these components are
transverse to the boundary and preserve generalized log canonicity, while
their ample class can be chosen large enough that \(D+A\) is nef. The
theorem supplies a terminating MMP starting on \(V\).

We verify the step types inductively.  Suppose \(V_i\) is
\(\Q\)-factorial generalized dlt.  By
\cite[Remark~2.15]{TsakanikasXie}, each birational step has the form
\[
 V_i\xrightarrow{g_i}Z_i\xleftarrow{h_i}V_{i+1},
\]
where \(g_i\) contracts a negative extremal ray and \(h_i\) is small and projective.
Suppose \(g_i\) is divisorial, and write \(R\) for its ray.  An exceptional
prime \(E\) has \(E\cdot R<0\): otherwise it would be effective,
exceptional, and \(g_i\)-nef, contrary to negativity.  There is only one
exceptional prime.  Indeed, a combination of two distinct exceptional
primes can be chosen to have degree zero on \(R\); negativity applied to
both signs of that combination would make it zero.

We show directly that \(Z_i\) is \(\Q\)-factorial.  For a prime divisor
\(T\) on \(Z_i\), let \(\widetilde T\) be its strict transform, choose a
curve \(C\) spanning \(R\), and put
\[
                 a=-\frac{\widetilde T\cdot C}{E\cdot C}\in\Q.
\]
Choose a positive integer \(m\) for which \(m(\widetilde T+aE)\) is
Cartier.  Its line has degree zero on \(R\), so
\cite[Theorem~1.5]{Xie2022} gives an actual Cartier line \(\mathcal L_T\)
on \(Z_i\) and an isomorphism
\[
 \OO_{V_i}\bigl(m(\widetilde T+aE)\bigr)\simeq g_i^*\mathcal L_T.
\]
On an open set trivializing \(\mathcal L_T\), the Cartier divisor
\(m(\widetilde T+aE)\) is principal upstairs.  Pushing this principal
divisor identity down through the birational \(g_i\) gives a principal
divisor \(mT\) on that open set.
Thus \(T\) is rational Cartier.  It follows that \(h_i\) is an isomorphism:
the pushforward of an \(h_i\)-ample Cartier divisor is rational Cartier
on \(Z_i\), and its pullback is the original divisor because \(h_i\)
is small.  It has degree zero on every contracted curve, so relative
ampleness rules out positive-dimensional fibers; a finite birational
map to the normal \(Z_i\) is an isomorphism.  If \(g_i\) is small,
the relatively ample canonical model \(h_i\) is the flip of
\cite[Definition~3.2]{HanLi}.  Thus the chosen birational steps are
divisorial contractions and flips.  The underlying variety of a
generalized dlt pair is klt, and these steps preserve the
\(\Q\)-factorial generalized dlt category
\cite[Remark~2.3 and Lemma~3.7]{HanLi}.  They introduce no prime
divisors on the new models. For completeness, on a common resolution
of one step \(V_i\dashrightarrow V_{i+1}\) carrying \(\boldsymbol M\),
taking the difference of the two generalized discrepancy formulas cancels
their common nef divisor. The resulting difference of the adjoint pullbacks is exceptional over the new
model and anti-nef over it. The negativity lemma makes this difference
effective. Composing these actual rational divisor comparisons gives
\(u^*D=v^*D_m+F\) with \(F\) effective and exceptional over the endpoint.
The endpoints in the two cases are the nef and Mori fiber endpoints supplied
by the terminating program theorem.
\end{proof}
